\documentclass[11pt,a4paper]{article}
\usepackage[margin=2.6cm]{geometry}
\usepackage{amsmath,amssymb,amsthm,mathtools,bm}
\usepackage{booktabs}
\usepackage{enumitem}
\usepackage[colorlinks=true,linkcolor=blue!50!black,citecolor=blue!50!black,urlcolor=blue!50!black]{hyperref}
\usepackage{microtype}
\usepackage{xcolor}

\newtheorem{proposition}{Proposition}
\newtheorem{lemma}{Lemma}
\theoremstyle{definition}
\newtheorem{definition}{Definition}
\newtheorem{exercise}{Exercise}
\theoremstyle{remark}
\newtheorem*{remark}{Remark}

\newcommand{\R}{\mathbb{R}}
\newcommand{\E}{\mathbb{E}}
\newcommand{\norm}[1]{\left\lVert #1 \right\rVert}
\newcommand{\ip}[2]{\left\langle #1, #2 \right\rangle}
\newcommand{\g}{\bm{g}}
\newcommand{\bmu}{\bm{\mu}}
\newcommand{\btheta}{\bm{\theta}}
\newcommand{\tr}{\operatorname{tr}}

\title{Is the auxiliary loss doing anything?\\[4pt]
\large Gradient geometry of multi-term objectives, with applications to\\
structural losses for precipitation super-resolution}
\author{Lecture notes for the \texttt{minkowski-loss} project}
\date{23 September 2026}

\begin{document}
\maketitle

\begin{abstract}
You train a network on $L_0 + \lambda L_k$ and want to know whether $L_k$ took any part in
shaping the result. This is harder than it looks. Loss \emph{values} carry no information
about influence. The obvious gradient diagnostic, comparing $\norm{\nabla L_0}$ with
$\lambda\norm{\nabla L_k}$, is only meaningful far from convergence. Adaptive optimisers
change the geometry in which ``how much'' is measured. These notes build the question up
from first principles:
\begin{itemize}[nosep]
\item one step of gradient descent;
\item the equilibrium of the combined objective, via linear-response theory;
\item stochastic dynamics, where the aux term's drift competes with minibatch noise;
\item Adam's preconditioned geometry;
\item statistical estimation from minibatches.
\end{itemize}
We then derive what the common weighting schemes (uncertainty weighting, GradNorm, MGDA,
PCGrad, black-box search) actually optimise. Finally, we apply everything to the five
structural losses of Study~1 and to the audit tool in \texttt{tools/gradient\_audit.py}.
We aim at understanding, not recipes. Each section ends with what it implies for the
measurement.
\end{abstract}

\tableofcontents
\bigskip

%=====================================================================
\section{The question, stated carefully}
%=====================================================================

Let $\btheta\in\R^D$ be the parameters ($D\approx 3.1\times10^7$ for our U-Net). We train on
\begin{equation}
  L(\btheta) = L_0(\btheta) + \lambda\,L_k(\btheta),
\end{equation}
where $L_0$ is the pixel MSE in log-normalised space, $L_k$ is one of the structural losses
(Minkowski, spectral, SSIM, wet area, optical flow), and $\lambda>0$ is fixed after a warm-up.
We write $\g_0=\nabla L_0$ and $\g_k=\nabla L_k$ for the full-data gradients. A single
minibatch $B$ gives unbiased noisy estimates $\hat\g_0,\hat\g_k$.

``Is $L_k$ active?'' has three non-equivalent meanings, and it pays to keep them apart:

\begin{enumerate}[label=(\alph*)]
\item \textbf{Instantaneous.} Does the term change the direction of the update
  \emph{now}? This is a local statement about $\g_0$ and $\lambda\g_k$ at one point.
\item \textbf{Cumulative.} Does the term change \emph{where training ends up}? This is a
  statement about the solution $\btheta_\lambda$ against $\btheta_0$. It depends on
  curvature and on noise, not just on gradient sizes.
\item \textbf{Observable.} Does the term change a metric we care about (MAE, exceedance
  ratio, the value of $L_k$ itself)? This is what the paper reports. It is a function of
  (b), seen through the metric.
\end{enumerate}

The audit numbers map onto these as follows:
\begin{itemize}[nosep]
\item the \emph{share} $r$ and the cosine address (a);
\item the ``own-loss'' test and the $\lambda$-scaling predictions of \S\ref{sec:equilibrium}
  address (b) and (c);
\item the MAE ``activity check'' addresses (c), but turns out to be a
  \emph{second-order} and hence weak test (\S\ref{sec:equilibrium}).
\end{itemize}

%=====================================================================
\section{Why loss values carry no information}
%=====================================================================

\begin{proposition}[Gauge freedom of a loss term]
For any $c>0$ and $a\in\R$, training on $L_0+\lambda L_k$ and on
$L_0+(\lambda/c)\,(c\,L_k+a)$ produces identical gradients, and therefore identical
trajectories under any first-order optimiser.
\end{proposition}
\begin{proof}
$\nabla[(\lambda/c)(cL_k+a)] = \lambda\nabla L_k$.
\end{proof}

The value of $L_k$ is thus defined only up to an affine map, and the value ratio
$\lambda L_k/L_0$ is not a physical quantity.
\begin{itemize}[nosep]
\item Our Minkowski loss takes values $\sim1$ and our spectral loss $\sim4\times10^{-3}$.
  That tells you nothing about which one steers more.
\item The additive constant matters more than it seems. SSIM has an \emph{unreachable
  floor} of $0.465$ on our data (\S\ref{sec:ssim}), so $93\%$ of its value is a constant with
  zero gradient.
\end{itemize}
Any weighting rule built on values inherits this arbitrariness. We will see in
\S\ref{sec:uw} that uncertainty weighting is exactly such a rule.

The invariant objects are the gradients $\lambda\g_k$, and more precisely their action
through the optimiser.

%=====================================================================
\section{One step of gradient descent}
%=====================================================================

Take one full-batch GD step $\Delta\btheta=-\eta(\g_0+\lambda\g_k)$. Let $\psi$ be the angle
between $\g_0$ and $\g_k$.

\begin{definition}[Share and parity weight]
\begin{equation}
  r \;=\; \frac{\lambda\norm{\g_k}}{\norm{\g_0}}, \qquad
  \lambda^* \;=\; \frac{\norm{\g_0}}{\norm{\g_k}}, \qquad\text{so}\quad r=\lambda/\lambda^*.
\end{equation}
\end{definition}

\paragraph{Direction.}
The update is deflected away from $-\g_0$ by an angle $\varphi$ with
\begin{equation}\label{eq:deflect}
  \tan\varphi = \frac{r\sin\psi}{1+r\cos\psi}.
\end{equation}
Two numbers from the first audit, spectral at $r=0.014$ and optical flow at $r=0.007$, give
$\varphi<0.8^\circ$. The update is, for all practical purposes, the MSE update.

\paragraph{What happens to each loss.}
To first order in $\eta$,
\begin{align}
  \Delta L_0 &= -\eta\left(\norm{\g_0}^2 + \lambda\ip{\g_0}{\g_k}\right)
             = -\eta\norm{\g_0}^2\,(1 + r\cos\psi), \label{eq:dL0}\\
  \Delta L_k &= -\eta\left(\lambda\norm{\g_k}^2 + \ip{\g_0}{\g_k}\right)
             = -\eta\norm{\g_0}\norm{\g_k}\,(r + \cos\psi). \label{eq:dLk}
\end{align}
Equation \eqref{eq:dLk} is the heart of the matter. The change in $L_k$ has a \emph{self}
term $r$ and a \emph{bystander} term $\cos\psi$ (the aux loss moves because the MSE moves).
The ratio self/bystander is $r/\cos\psi$. Three regimes follow:
\begin{itemize}[nosep]
\item $r\ll|\cos\psi|$: $L_k$ is a \emph{bystander}. Whatever it does is inherited from the
  MSE. A model trained this way should score on $L_k$ like a vanilla model trained for as
  long.
\item $r\gg|\cos\psi|$: $L_k$ drives its own descent.
\item In high dimension, independent directions have $\cos\psi = O(D^{-1/2})$. A small
  measured $\cos\psi$ (Minkowski: $-0.02$ in the first audit) means that even a modest $r$
  is enough for the self term to dominate. \emph{Small $r$ is not the same as inactive;
  what matters is $r$ relative to $|\cos\psi|$.}
\end{itemize}
Equation \eqref{eq:dL0} shows the first-order cost to $L_0$: only a \emph{conflicting}
aux gradient ($\cos\psi<0$) slows the MSE. An orthogonal one costs nothing at first order.
Its cost appears at second order, which is the subject of the next section.

\begin{remark}
$\lambda^*$ is a \emph{unit}, not an optimum. It is the weight at which the two terms have
equal say in the instantaneous direction. Our Minkowski loss works at $r\approx 2$--$5$ and
games the metrics at $10\times$ its weight. Nothing forces the good $\lambda$ to be near
$\lambda^*$. Its use is to express every loss's weight in the same dimensionless unit, which
is what a fair comparison needs.
\end{remark}

%=====================================================================
\section{Where training ends up: equilibrium and linear response}\label{sec:equilibrium}
%=====================================================================

One step tells you about the direction. The solution is determined by a balance of forces.
Physicists will recognise the structure: the auxiliary loss is an external field, the base
loss is the potential, and the Hessian is the stiffness.

\subsection{Perturbation theory around the base minimum}

Let $\btheta_0$ be a non-degenerate local minimum of $L_0$, with Hessian
$H=\nabla^2L_0(\btheta_0)\succ0$. Write $\g_k=\nabla L_k(\btheta_0)$. By the implicit function
theorem the minimiser of $L_0+\lambda L_k$ satisfies
\begin{equation}\label{eq:response}
  \btheta_\lambda = \btheta_0 - \lambda H^{-1}\g_k + O(\lambda^2).
\end{equation}

\begin{proposition}[Response of the two losses]\label{prop:response}
Define the \emph{susceptibility} $\chi_k = \g_k^\top H^{-1}\g_k > 0$. Then
\begin{align}
  L_k(\btheta_\lambda)-L_k(\btheta_0) &= -\lambda\,\chi_k + O(\lambda^2), \label{eq:lin}\\
  L_0(\btheta_\lambda)-L_0(\btheta_0) &= \tfrac12\lambda^2\,\chi_k + O(\lambda^3). \label{eq:quad}
\end{align}
Eliminating $\lambda$, the local Pareto front is a parabola:
\begin{equation}\label{eq:pareto}
  \Delta L_0 \approx \frac{(\Delta L_k)^2}{2\chi_k}.
\end{equation}
\end{proposition}
\begin{proof}
Taylor-expand both losses around $\btheta_0$ using \eqref{eq:response}. For $L_k$ the linear
term is $\g_k^\top\Delta\btheta = -\lambda\chi_k$. For $L_0$ the gradient vanishes at
$\btheta_0$, leaving $\tfrac12\Delta\btheta^\top H\Delta\btheta=\tfrac12\lambda^2\chi_k$.
\end{proof}

Three consequences shape the whole audit.

\paragraph{1. The MAE activity check is a second-order test.}
The cost to the base loss grows like $\lambda^2$, and the gain on the aux loss like
$\lambda$. So at small $\lambda$ a genuinely active term can improve its own objective
noticeably while moving MSE or MAE imperceptibly. ``$\Delta\mathrm{MAE}<1\%$ means
inactive'' is therefore a sufficient signal only when combined with ``and its own loss did
not improve either''. The \textbf{own-loss test} (the trained model evaluated on its own
$L_k$, against vanilla) is first order and strictly more sensitive. It is the primary
evidence in the audit.

\paragraph{2. The right norm is $H^{-1}$, not Euclidean.}
$\chi_k$ weights the aux gradient by the \emph{inverse} stiffness of the base loss.
\begin{itemize}[nosep]
\item Directions in which the MSE is flat (small curvature) respond enormously to a small
  push. Directions in which it is stiff barely move.
\item The Euclidean share $r$ ignores this completely, so it is a proxy for (b) at best.
\item This is also where reward hacking lives. The grid-aligned filaments of the
  $w=10^{-3}$ Minkowski run (\texttt{DECISIONS.md} \S2) are, almost by definition, a
  direction the MSE does not resist much and the Minkowski loss rewards.
\end{itemize}

\paragraph{3. At equilibrium the full gradients balance exactly.}
At a stationary point of $L_0+\lambda L_k$ one has $\g_0 = -\lambda\g_k$. So
\begin{equation}
  r = 1, \qquad \cos\psi = -1 \qquad\text{(full batch, at the combined optimum).}
\end{equation}
This gives a clean diagnostic for a model trained \emph{with} $L_k$: at its own checkpoint
the noise-corrected share should approach $1$ and the cosine $-1$. The size of the deviation
measures how far from equilibrium training stopped.
\begin{itemize}[nosep]
\item Conversely, at the \emph{vanilla} checkpoint $\g_0\approx0$ (it is a minimum of $L_0$).
  So $\lambda^*=\norm{\g_0}/\norm{\g_k}\to0$ and $r\to\infty$ for \emph{every} aux loss.
\item The parity weight evaluated at the ERM solution is therefore not a well-posed
  quantity, contrary to the prescription in \texttt{DECISIONS.md} \S9 and in Chang \& Sapsis.
  What one measures there, with minibatches, is mostly the noise scale of $\hat\g_0$
  (\S\ref{sec:noise}).
\item The well-posed places to measure parity are the \textbf{initialisation}, or early
  training, where both gradients are signal-dominated. There the ratio governs the
  transient that decides which basin you enter.
\end{itemize}

\subsection{The overparameterised caveat}
With $D\gg$ (effective number of constraints), $H$ is singular. The minimisers of $L_0$
form a manifold $\mathcal M$, and \eqref{eq:response} holds only in the directions transverse
to $\mathcal M$.
\begin{itemize}[nosep]
\item \emph{Along} $\mathcal M$, $L_0$ exerts no restoring force. An arbitrarily small
  $\lambda$ then drives the solution, given enough time, to $\arg\min_{\mathcal M}L_k$: the
  aux loss acts as a tie-breaker at zero cost to $L_0$.
\item In that regime $\chi_k=\infty$ in the flat directions. MAE-based checks see nothing,
  while $L_k$ can move by $O(1)$.
\item Our U-Net does not interpolate the data (validation MSE is far from zero, and the
  targets are noisy), and training is finite. So we are in an intermediate regime: some
  flat directions, many stiff ones. It is one more reason to trust the own-loss test over
  the MAE check.
\end{itemize}

%=====================================================================
\section{Noise: drift against diffusion}\label{sec:noise}
%=====================================================================

Minibatch gradients are $\hat\g=\bmu+\bm\xi$ with $\E\bm\xi=0$ and $\operatorname{Cov}\bm\xi=\Sigma/B$
(approximately, for batch size $B$). The aux term contributes a \emph{drift}
$\lambda\bmu_k$. The base loss contributes a restoring force plus diffusion. Whether the drift
matters is a signal-to-noise question.

\subsection{A one-dimensional model}
Project the dynamics onto the unit vector $\bm e=\bmu_k/\norm{\bmu_k}$. Near a minimum of $L_0$
with curvature $h=\bm e^\top H\bm e$ along $\bm e$, SGD with learning rate $\eta$ is a discretised
Ornstein--Uhlenbeck process:
\begin{equation}
  x_{t+1} = x_t - \eta\,(h\,x_t + \lambda\norm{\bmu_k}) + \eta\,\sigma_e\,\zeta_t,
  \qquad \sigma_e^2=\bm e^\top\Sigma\,\bm e/B,\quad \zeta_t\sim\mathcal N(0,1).
\end{equation}
For small $\eta h$ the stationary distribution has
\begin{equation}
  \text{mean shift}\;\; m=-\frac{\lambda\norm{\bmu_k}}{h},\qquad
  \text{spread}\;\; s=\sigma_e\sqrt{\frac{\eta}{2h}},\qquad
  \frac{|m|}{s}=\frac{\lambda\norm{\bmu_k}}{\sigma_e}\sqrt{\frac{2}{\eta h}}.
\end{equation}
The relaxation time is $1/(\eta h)$ steps. Three lessons:
\begin{itemize}[nosep]
\item An aux term can shift the \emph{mean} solution while being invisible in any single
  step, because $\lambda\norm{\bmu_k}\ll\sigma_e$. A consistent drift accumulates; noise
  does not.
\item Without restoring force ($h\to0$, a flat direction) the drift wins after
  $T^*\approx\big(\sigma_e/(\lambda\norm{\bmu_k})\big)^2$ steps.
\item Learning-rate decay ($\eta\downarrow$, as \texttt{ReduceLROnPlateau} does) \emph{raises}
  $|m|/s$. Late training is where small, consistent aux terms leave their mark.
\end{itemize}

\subsection{What a per-batch norm measures}
\begin{lemma}\label{lem:noise}
$\E\norm{\hat\g}^2=\norm{\bmu}^2+\tr\Sigma/B$.
\end{lemma}
Near a minimum of $L_0$, $\norm{\bmu_0}\to0$ while $\tr\Sigma_0/B$ does not. A per-batch
$\norm{\hat\g_0}$ is then a measure of \emph{noise}, and a parity ratio built from per-batch
norms compares the aux signal with the MSE noise. That is a meaningful comparison (it is
$\lambda\norm{\bmu_k}/\sigma$ above, up to the $\sqrt{\eta h}$ factor), but it is not the
parity of \S3.

The first audit (\texttt{tools/gradient\_audit\_v0}) used per-batch norms at the vanilla
checkpoint. Its ``$\lambda^*$'' values are therefore noise-referenced. The qualitative
ordering is still informative, but the ``$70\times$'' and ``$140\times$'' factors should be
read in that light.

\subsection{Unbiased estimation}
With $N$ minibatch gradients $\hat\g^{(i)}$, mean $\bar\g$ and mean squared norm
$S=\frac1N\sum_i\norm{\hat\g^{(i)}}^2$:
\begin{equation}\label{eq:unbiased}
  \widehat{\norm{\bmu}^2} = \frac{N\norm{\bar\g}^2-S}{N-1},\qquad
  \widehat{\ip{\bmu_0}{\bmu_k}} = \frac{N\ip{\bar\g_0}{\bar\g_k}-\frac1N\sum_i\ip{\hat\g_0^{(i)}}{\hat\g_k^{(i)}}}{N-1}.
\end{equation}
Both follow from Lemma~\ref{lem:noise} and its bilinear analogue. The second one matters
because $\hat\g_0$ and $\hat\g_k$ are computed on the \emph{same} batch, so their noises are
correlated. The naive inner product of mean gradients is biased by that cross-covariance,
by $\tr\Sigma_{0k}/(NB)$. So even for orthogonal signals the naive cosine is not centred on
zero when $N$ is small.

The tool reports:
\begin{itemize}[nosep]
\item the noise-corrected norms and cosine;
\item the \emph{noise fraction} $1-\widehat{\norm{\bmu}^2}/\text{median}\norm{\hat\g}^2$;
\item jackknife standard errors over $K$ groups of batches.
\end{itemize}
When $\widehat{\norm{\bmu_0}^2}\le0$, the MSE signal is statistically indistinguishable from
zero, and the tool says so instead of printing a meaningless ratio.

%=====================================================================
\section{Adam changes the geometry}\label{sec:adam}
%=====================================================================

We train with AdamW. Its update is
\begin{equation}
  \btheta \leftarrow \btheta(1-\eta w) - \eta\,\frac{\hat{\bm m}}{\sqrt{\hat{\bm v}}+\epsilon},
\end{equation}
with $\hat{\bm m}\approx\E\hat\g$ and $\hat{\bm v}\approx\E\hat\g^{\odot2}$ (coordinate-wise), and
decoupled weight decay $w$.

\paragraph{Preconditioned share.}
Holding $\hat{\bm v}$ fixed, adding $\lambda\g_k$ changes the update by
$-\eta\lambda P\g_k$ with $P=\operatorname{diag}\big(1/(\sqrt{\hat{\bm v}}+\epsilon)\big)$. The relevant
share is therefore
\begin{equation}
  r_P = \frac{\lambda\norm{P\g_k}}{\norm{P\g_0}}.
\end{equation}
The tool computes it with $\hat{\bm v}$ taken from each checkpoint's saved optimizer state.
\begin{itemize}[nosep]
\item Near convergence $\hat{\bm v}$ is dominated by noise, so $P\approx\operatorname{diag}(1/\sigma_i)$.
  Adam then measures each coordinate's gradient \emph{in units of its own noise}, and $r_P$
  is a signal-to-noise-weighted share.
\item A loss whose gradient concentrates in coordinates where the MSE is quiet gets
  amplified. One spread over noisy coordinates gets suppressed.
\item In the smoke test (4 batches, not to be quoted) the preconditioner cut SSIM's share
  by almost an order of magnitude and wet area's by seven. Euclidean and preconditioned
  shares can tell different stories.
\end{itemize}

\paragraph{Invariances.}
Adam is invariant to a global rescaling of the \emph{total} loss (up to $\epsilon$). It is
\emph{not} invariant to $\lambda$, since that changes the relative composition of $\hat{\bm m}$
and $\hat{\bm v}$. So ``Adam normalises the gradients, so weights don't matter'' is wrong.
They matter through the direction.

\paragraph{The AdamW fixed point is preconditioner-dependent.}
Stationarity requires $P\nabla L=-w\btheta$, not $\nabla L=0$. The linear response
\eqref{eq:response} becomes
\begin{equation}
  \btheta_\lambda-\btheta_0 \approx -\lambda\,(PH+wI)^{-1}P\,\g_k \quad(\text{$P$ held fixed}),
\end{equation}
which reduces to $-\lambda H^{-1}\g_k$ as $w\to0$. The equilibrium itself therefore depends
on the optimiser's statistics. This is one more reason the Euclidean parity rule is only a
first approximation.

\paragraph{Clipping and mixed precision.}
\begin{itemize}[nosep]
\item \emph{Clipping.} Global-norm clipping at $c=1$ rescales the total update. It changes
  the step size, not the share, and only when $\norm{\hat\g}>1$. The tool reports how often
  that happens.
\item \emph{fp16 underflow.} Training runs in fp16 autocast with a \texttt{GradScaler}. The
  scaled backward pass multiplies gradients by $S\sim2^{16}$ before unscaling. A very small
  $\lambda$ times small per-element aux gradients can nevertheless underflow in fp16
  intermediates (the smallest subnormal is $\approx6\times10^{-8}$).
  \textbf{This is a hypothesis, not a finding.} But for $\lambda=4\times10^{-5}$ (optical
  flow) it is worth checking, because it would make a small weight even smaller in
  practice. The audit runs in fp32 and cannot see it. A direct test compares one fp16 and
  one fp32 backward on the same batch.
\end{itemize}

%=====================================================================
\section{What the standard weighting schemes optimise}
%=====================================================================

\subsection{Uncertainty (homoscedastic) weighting}\label{sec:uw}
Kendall, Gal \& Cipolla (2018) minimise
\begin{equation}
  \mathcal L(\btheta,\bm\sigma)=\sum_k\frac{1}{2\sigma_k^2}L_k(\btheta)+\log\sigma_k.
\end{equation}
\begin{proposition}
At the optimum over $\sigma_k$ one has $\sigma_k^2=L_k$. The effective weight is then
$1/(2L_k)$, and the $\btheta$-gradient becomes
\begin{equation}
  \nabla_{\btheta}\mathcal L=\sum_k\frac{\nabla L_k}{2L_k}=\tfrac12\nabla\sum_k\log L_k.
\end{equation}
\end{proposition}
\begin{proof}
$\partial_{\sigma_k}\mathcal L=-L_k/\sigma_k^3+1/\sigma_k=0$.
\end{proof}
So uncertainty weighting equalises \emph{relative} (logarithmic) progress across terms.
\begin{itemize}[nosep]
\item It is invariant to rescaling $L_k\to cL_k$, which is good.
\item It is \emph{not} invariant to offsets. With an unreachable floor $L_k=f+\ell_k$, the
  weight is $1/(2(f+\ell_k))$. For SSIM, $f=0.465$ dominates, so the floor sets the weight.
\item It is a Gaussian-likelihood argument. None of our structural losses is a Gaussian
  negative log-likelihood, so the $\sigma_k$ have no probabilistic meaning here.
\item The weight is \emph{learned} per loss. In a comparison of losses, that means each
  loss gets an uncontrolled, loss-specific $\lambda$, which is exactly what a fair comparison
  must avoid.
\end{itemize}
It is well suited to balancing the three Minkowski channels (area, perimeter, $\chi$)
against each other, since they are commensurable; that is where the project already uses it.

\subsection{GradNorm}
Chen et al.\ (2018) adapt $\lambda_k(t)$ so that the weighted gradient norms (on the last
shared layer) track a common target, modulated by each task's relative training rate. It is
\emph{dynamic parity}: it drives $r\to\text{target}$ throughout training. It solves the
units problem. But it fixes the operating point at parity, whereas our experience (the
Minkowski loss works at $r\approx2$--$5$ and hacks at $\approx10\times$ that weight) says the
right point is loss-specific.

\subsection{MGDA and PCGrad}
\begin{itemize}[nosep]
\item \emph{MGDA} (Sener \& Koltun 2018) picks the min-norm point of the convex hull of
  $\{\g_k\}$, which gives a common descent direction and converges to Pareto stationarity.
  It encodes no preference between objectives, and here we have a strong one: the MSE is
  primary.
\item \emph{PCGrad} (Yu et al.\ 2020) projects out conflicting components when
  $\cos\psi<0$. With $\cos\psi\approx0$, as measured for Minkowski, it does essentially
  nothing.
\end{itemize}

\subsection{Black-box search (Optuna, Bayesian optimisation)}
Treat $\lambda$ as a hyperparameter and maximise a validation score. This is principled,
with two caveats:
\begin{itemize}[nosep]
\item \emph{The result is the chosen metric.} Choosing it is the actual scientific decision,
  and it must not be any loss's own objective.
\item \emph{The search is one-dimensional per loss and expensive.} A handful of bracketed
  runs covers it. Search pays off for joint spaces ($\lambda\times$ anneal $\times$ tail
  sampling), ideally on short proxy runs with pruning and an equal budget per loss.
\end{itemize}

\subsection{What a fair comparison needs}
Put every loss in the same dimensionless unit, then probe the response.
\begin{enumerate}[nosep]
\item Measure the parity unit $\lambda^*$ where it is well posed: at initialisation or
  early training (\S\ref{sec:equilibrium}), Euclidean and Adam-preconditioned.
\item Train a fixed log-bracket around it, reaching above parity.
\item Evaluate at a fixed budget, on the last checkpoint.
\item Check the \emph{signatures} of Proposition~\ref{prop:response}:
  \begin{itemize}[nosep]
  \item own-loss gain $\propto\lambda$;
  \item base-loss cost $\propto\lambda^2$;
  \item at the trained checkpoint, share $\to1$ and cosine $\to-1$.
  \end{itemize}
  A loss that shows no linear own-loss gain across the bracket is genuinely inert, not
  merely under-weighted.
\end{enumerate}
This is the protocol in \texttt{DECISIONS.md} \S16, with the correction that parity is
measured at initialisation rather than at the ERM solution.

%=====================================================================
\section{Case studies: where each structural loss puts its gradient}
%=====================================================================

The gradient of a loss is a field over pixels, $\partial L_k/\partial\hat y_p$, pulled back
through the network. Its \emph{support} (which pixels receive signal) explains most of the
audit.

\subsection{Minkowski: threshold-localised gradients and the pixel-mass bottleneck}
The analytical loss replaces the indicator $\mathbb 1[\hat y_p\ge u]$ by
$s_p(u)=\sigma\!\big((\hat y_p-u)/\tau_u\big)$, with $\tau_u=\max(0.1\,u,\tau_{\min})$.
Area, perimeter and Euler characteristic are smooth functionals of $\{s_p(u)\}$. Since
$\sigma'(z)=\sigma(z)(1-\sigma(z))$ decays like $e^{-|z|}$, pixel $p$ receives gradient from
threshold $u$ only if $|\hat y_p-u|\lesssim$ a few $\tau_u$. The number of pixels in that
band is
\begin{equation}
  N_u\;\approx\;|D|\int_{u-c\tau_u}^{u+c\tau_u} f(y)\,dy\;\approx\;2c\,|D|\,\tau_u f(u)\;\propto\; u\,f(u),
\end{equation}
where $f$ is the pixel-intensity density.
\begin{itemize}[nosep]
\item For a heavy tail $f(u)\propto u^{-\alpha}$ with $\alpha>1$, $N_u\propto u^{1-\alpha}$
  decreases with $u$. The thresholds at $31$--$150$\,mm/h see almost no pixels in their band.
\item That is the measured result that 83\% of the gradient comes from $u\le0.28$\,mm/h and
  $\approx0$ from $u\ge31$ (\texttt{DECISIONS.md} \S5), derived rather than observed.
\item Widening $\tau$ (annealing) raises every band proportionally, so the \emph{share}
  cannot change. That is the anneal-sweep result.
\item The remedies that change $N_u$ itself are tail-sample selection and threshold
  reweighting.
\end{itemize}

\subsection{SSIM: an unreachable floor}\label{sec:ssim}
Per window, $\mathrm{SSIM}=\dfrac{(2\mu_x\mu_y+c_1)(2\sigma_{xy}+c_2)}{(\mu_x^2+\mu_y^2+c_1)(\sigma_x^2+\sigma_y^2+c_2)+\varepsilon}$,
with $c_{1,2}=(k_{1,2}R)^2$, where $R$ is the per-sample data range. The implementation also
adds $\varepsilon=10^{-6}$ to the denominator.

In an exactly dry window of both prediction and target, $\mu=\sigma=0$, so
$\mathrm{SSIM}=c_1c_2/(c_1c_2+\varepsilon)$. With $k_1=0.01$, $k_2=0.03$ and $R\le1$ in
normalised units, $c_1c_2\le9\times10^{-8}R^4\ll\varepsilon$. Hence $\mathrm{SSIM}\approx0$
and the per-window loss $(1-\mathrm{SSIM})/2\approx\tfrac12$, \emph{even for a perfect
prediction}. The dataset-level floor is
\begin{equation}
  L_{\text{SSIM}}^{\min}\approx\tfrac12\times(\text{fraction of dry windows})
  = 0.465 \;\Rightarrow\; \approx93\%\text{ of windows dry}.
\end{equation}
The gradient lives in the remaining $\approx7\%$, over a dynamic range of $0.465\to0.5$. Two
fixes:
\begin{itemize}[nosep]
\item a fixed data range ($R=1$ in normalised space) with $\varepsilon\ll c_1c_2$, which
  makes dry-on-dry windows score $1$;
\item or restricting the mean to windows that are wet in either field.
\end{itemize}

\subsection{Spectral, wet area, optical flow}
\begin{itemize}[nosep]
\item \textbf{Spectral} (FFT-magnitude energy score) is translation-invariant, so its
  gradient is orthogonal to every translation of the prediction. It cannot penalise
  displacement. Its value is an $L^1$ over all rfft bins, dominated by high-wavenumber bins
  of nearly dry fields, which explains its tiny scale.
\item \textbf{Wet area} is the single-threshold ($u=0.1$\,mm/h, $\tau=0.05$) special case of
  the Minkowski area channel. Its gradient support is the drizzle band only, so it is blind
  to intensity. It can be ``active'' ($r\gg1$) and still useless for extremes.
\item \textbf{Optical flow} uses the coarse input as its teacher. That teacher is smooth
  with suppressed peaks, so the gradient points toward resembling the input.
\end{itemize}

%=====================================================================
\section{Reading the audit}
%=====================================================================

\texttt{tools/gradient\_audit.py} reports, per checkpoint and loss:

\begin{center}
\begin{tabular}{@{}lll@{}}
\toprule
Column & Definition & Answers \\ \midrule
parity $\lambda^*$ & $\widehat{\norm{\bmu_0}}/\widehat{\norm{\bmu_k}}$, eq.~\eqref{eq:unbiased} & the unit of weight \\
share $r$ & $\lambda/\lambda^*$ & instantaneous influence (a) \\
per-batch $\lambda^*$ & median of $\norm{\hat\g_0}/\norm{\hat\g_k}$ & noise-referenced (the v0 number) \\
cos, cos corr. & naive and noise-corrected & bystander vs self, eq.~\eqref{eq:dLk} \\
Adam share & $\lambda\norm{P\bar\g_k}/\norm{P\bar\g_0}$ & influence on the real update \\
noise fraction & $1-\widehat{\norm{\bmu}^2}/\mathrm{med}\norm{\hat\g}^2$ & is this norm signal? \\
clip engaged & share of batches with $\norm{\hat\g_0+\lambda\hat\g_k}>1$ & does clipping act? \\
own vs vanilla & paired, bootstrap CI & first-order effect (b,c) \\ \bottomrule
\end{tabular}
\end{center}

\noindent Checklist:
\begin{enumerate}[nosep]
\item \textbf{Initialisation:} is $r$ (Euclidean) small, and is it small relative to
  $|\cos\psi|$? That tells you whether the term could have shaped the early trajectory.
\item \textbf{Vanilla checkpoint:} expect $\widehat{\norm{\bmu_0}}\approx0$ and a
  noise fraction near $1$ for the MSE. Parity is undefined there.
\item \textbf{Own checkpoint of a run trained with $L_k$:} does the share approach $1$ and
  the cosine $-1$? If not, the run was stopped far from the combined equilibrium, or the term
  never mattered.
\item \textbf{Own-loss test:} does the trained model beat vanilla on its own $L_k$, with
  a CI that excludes 0? This is the decisive first-order evidence.
\item \textbf{Across a bracket} (after the M4 reruns): is the own-loss gain linear in
  $\lambda$ and the MSE cost quadratic? That confirms the linear-response regime and locates
  where it breaks down (hacking).
\end{enumerate}

%=====================================================================
\section{Exercises}
%=====================================================================

\begin{exercise}
Show that $\tan\varphi=r\sin\psi/(1+r\cos\psi)$ in \eqref{eq:deflect}. For fixed
$\psi>\pi/2$, how does $\varphi$ behave as $r$ grows, and at which $r$ does the update stop
decreasing $L_0$? What happens when $\cos\psi=-1$?
\end{exercise}

\begin{exercise}
Prove Lemma~\ref{lem:noise} and derive the unbiased estimators in \eqref{eq:unbiased}.
Why does the naive cosine of two mean gradients computed on the same batches not tend to
$0$ for independent signals when $N$ is small?
\end{exercise}

\begin{exercise}
Using Proposition~\ref{prop:response}, show that if a run at $\lambda$ improves its own loss
by $\delta$ at an MSE cost $\epsilon$, then a run at $\lambda/3$ should improve by
$\approx\delta/3$ at a cost of $\approx\epsilon/9$. How would you use this to tell
``inert'' from ``under-weighted'' with three runs?
\end{exercise}

\begin{exercise}
In the OU model of \S\ref{sec:noise}, compute the number of steps after which the mean
shift reaches $1-e^{-1}$ of its stationary value. With $\eta=5\times10^{-5}$, how large
must $h$ be for a 25-epoch run (about $4.5\times10^5$ steps at batch 128 on 2.3\,M patches)
to equilibrate?
\end{exercise}

\begin{exercise}
Show that uncertainty weighting applied to $L_k=f+\ell_k$ with a constant floor $f>0$ gives
effective weight $1/(2(f+\ell_k))$. For SSIM ($f=0.465$, $\ell_k\in[0,0.035]$), how much does
the effective weight change over the whole achievable range of $\ell_k$? What does this say
about what the learned $\sigma_k$ can encode?
\end{exercise}

\subsection*{Brief solutions}
\begin{enumerate}[label=\arabic*.,nosep]
\item Decompose $\lambda\g_k$ into components parallel and perpendicular to $\g_0$. For
  $\psi>\pi/2$, $\varphi$ increases monotonically from $0$ towards $\psi$ as $r\to\infty$. It
  passes $90^\circ$ at $r=-1/\cos\psi$, where the update has no component along $-\g_0$ and
  $L_0$ stops decreasing (eq.~\eqref{eq:dL0}). At $\cos\psi=-1$ and $r=1$ the two gradients
  cancel exactly: that is the equilibrium of \S\ref{sec:equilibrium}.
\item Expand $\E\norm{\bar\g}^2=\norm{\bmu}^2+\tr\Sigma/(NB)$ and eliminate $\tr\Sigma$
  using $\E S$. Same-batch noise is correlated with covariance $\Sigma_{0k}$, which biases
  $\ip{\bar\g_0}{\bar\g_k}$ by $\tr\Sigma_{0k}/(NB)$.
\item Linear gain and quadratic cost. Three bracketed runs with no linear own-loss trend
  mean $\chi_k\approx0$ (inert). A linear trend that is merely small means under-weighted.
\item $t\approx1/(\eta h)$. Equilibration within $4.5\times10^5$ steps needs
  $h\gtrsim 1/(\eta\cdot4.5\times10^5)\approx 0.044$ along $\bm e$. Flatter directions
  do not equilibrate within the run.
\item The weight varies by $<8\%$ over the achievable range, so $\sigma_k$ encodes
  essentially only the floor.
\end{enumerate}


%=====================================================================
\appendix
\section{Measurements, 23 September 2026}\label{app:results}
%=====================================================================

\texttt{tools/gradient\_audit.py} with $N=512$ batches of $128$ patches ($65{,}536$ patches,
drawn with the training sampler), $K=8$ jackknife groups, anneal $0.1$, fp32, GPU~1. Raw
output: \texttt{eval\_results/gradient\_audit/2026-09-23/}. Weights $\lambda$ are those the
Study-1 runs were trained with.

\subsection*{A.1 Parity at initialisation (where it is well posed)}
At initialisation the MSE gradient is almost pure signal (noise fraction $0.01$), so
$\lambda^*$ is well defined.
\begin{center}
\begin{tabular}{@{}lrrrr@{}}
\toprule
loss & $\lambda$ used & $\lambda^*$ & $r=\lambda/\lambda^*$ & $\cos\psi$ \\ \midrule
Minkowski   & $1\times10^{-4}$   & $1.05\times10^{-3}$ & $0.095$              & $+0.78$ \\
spectral    & $5\times10^{-4}$   & $0.298$             & $1.7\times10^{-3}$   & $+0.96$ \\
SSIM        & $1.6\times10^{-4}$ & $0.932$             & $1.7\times10^{-4}$   & $-0.77$ \\
wet area    & $2\times10^{-2}$   & $2.00\times10^{-2}$ & $1.0$                & $+0.85$ \\
optical flow& $4\times10^{-5}$   & $1.09$              & $3.7\times10^{-5}$   & $+0.79$ \\ \bottomrule
\end{tabular}
\end{center}
At initialisation every aux gradient is strongly aligned with the MSE gradient
($|\cos\psi|\approx0.8$--$0.96$), so every term starts in the bystander regime
$r\ll|\cos\psi|$ of eq.~\eqref{eq:dLk}. Even Minkowski starts $10\times$ below parity, and the
term only takes over later, as the MSE gradient shrinks. Wet area is at parity and clips
$64\%$ of initial batches.

\subsection*{A.2 The vanilla checkpoint: parity is ill-posed there}
At the converged vanilla checkpoint the MSE noise fraction is $0.96$: the per-batch norm
($3.0\times10^{-4}$) is almost all noise around a signal of $6.1\times10^{-5}$. The per-batch
parity values of the first audit are reproduced (spectral $3.9\times10^{-2}$, v0:
$3.6\times10^{-2}$). The noise-corrected values are $4$--$5\times$ smaller (spectral
$9.2\times10^{-3}$). This is Lemma~\ref{lem:noise} in action. The ``$70\times$ below parity''
of the first audit is $\approx18\times$ after noise correction, and $\approx600\times$ at
initialisation. The factor depends entirely on where it is measured.

\subsection*{A.3 The equilibrium test at each run's own checkpoint}
For a term that trained actively, \S\ref{sec:equilibrium} predicts $r\to1$ and $\cos\psi\to-1$.
\begin{center}
\begin{tabular}{@{}lrrrl@{}}
\toprule
run (own loss) & $r$ & $\cos\psi$ (corr.) & Adam share & verdict \\ \midrule
Minkowski   & $1.22$   & $-0.986$ & $1.10$   & at equilibrium: active \\
wet area    & $0.74$   & $-0.967$ & $1.03$   & at equilibrium: active \\
SSIM        & $0.157$  & $+0.470$ & $0.038$  & far from it: bystander \\
spectral    & $0.021$  & $+0.915$ & $0.022$  & far from it: bystander \\
optical flow& $0.0066$ & $-0.474$ & $0.008$  & far from it: inert \\ \bottomrule
\end{tabular}
\end{center}
The two active runs sit almost exactly on the equilibrium line, with Euclidean and Adam
shares agreeing. The three others have shares $\ll1$ and positive or weakly negative cosines.
Their own gradient never balanced the MSE, and in two of them it still points the same way.

\subsection*{A.4 The own-loss test}
Paired over the same $512$ batches, with bootstrap 95\% intervals, relative to vanilla:
Minkowski $-64.8\%$ $[-65.9,-63.8]$, wet area $-13.9\%$ $[-15.8,-12.1]$, optical flow
$-2.6\%$ $[-6.4,+1.2]$, spectral $+2.5\%$ $[+0.9,+4.4]$, SSIM $+0.1\%$ $[+0.0,+0.2]$.
\begin{itemize}[nosep]
\item Spectral and SSIM are \emph{worse} on their own objective than the vanilla model.
\item The 34-epoch vanilla model is also worse than the 64-epoch one on spectral
  ($4.18$ vs $4.10\times10^{-3}$). So the sign is explained by under-training at 25
  epochs, not by the term.
\item Both tests agree. Spectral, SSIM and optical flow were bystanders. Minkowski and wet
  area were active, and wet area's activity is confined to the drizzle band, as its
  gradient support predicts.
\end{itemize}

\subsection*{A.5 Consequences for the reruns}
\begin{itemize}[nosep]
\item Use the \emph{initialisation} $\lambda^*$ as the unit. By it, the Study-1 weights were
  $\sim10^{-1}$ (Minkowski), $\sim10^{-3}$ (spectral), $\sim10^{-4}$ (SSIM),
  $\sim4\times10^{-5}$ (optical flow) and $1$ (wet area) times parity. Minkowski became
  active only because $\norm{\g_0}$ collapsed during training (A.3). The competing losses
  would need orders of magnitude more weight to follow the same route.
\item A bracket expressed in init-parity units cannot be copied from Minkowski's working
  point. Minkowski was active at $0.1\lambda^*_{\text{init}}$, yet that same fraction leaves
  the others inert, because their gradients are dense and aligned with the MSE.
  \textbf{The bracket should instead be placed so that the equilibrium test (A.3) is
  satisfied at the end of training}, and checked with the own-loss test.
\end{itemize}

%=====================================================================
\section*{References}
%=====================================================================
\begin{itemize}[leftmargin=*,nosep]
\item Z.~Chen, V.~Badrinarayanan, C.-Y.~Lee, A.~Rabinovich. GradNorm: gradient
  normalization for adaptive loss balancing in deep multitask networks. \emph{ICML} 2018.
\item A.~Kendall, Y.~Gal, R.~Cipolla. Multi-task learning using uncertainty to weigh losses
  for scene geometry and semantics. \emph{CVPR} 2018.
\item O.~Sener, V.~Koltun. Multi-task learning as multi-objective optimization.
  \emph{NeurIPS} 2018.
\item T.~Yu, S.~Kumar, A.~Gupta, S.~Levine, K.~Hausman, C.~Finn. Gradient surgery for
  multi-task learning. \emph{NeurIPS} 2020.
\item D.~P.~Kingma, J.~Ba. Adam: a method for stochastic optimization. \emph{ICLR} 2015.
\item I.~Loshchilov, F.~Hutter. Decoupled weight decay regularization. \emph{ICLR} 2019.
\item S.~Mandt, M.~D.~Hoffman, D.~M.~Blei. Stochastic gradient descent as approximate
  Bayesian inference. \emph{JMLR} 18, 2017.
\item Z.~Wang, A.~C.~Bovik, H.~R.~Sheikh, E.~P.~Simoncelli. Image quality assessment: from
  error visibility to structural similarity. \emph{IEEE Trans.\ Image Process.} 13, 2004.
\item Chang \& Sapsis (2026). Extreme Event Aware Learning. doi:10.1038/s41467-026-76811-x.
\end{itemize}
\medskip
\noindent\small Author lists and venues of the machine-learning references are from memory.
Check them before citing in the paper.

\end{document}
